% !TeX root = ../main.tex

\chapter{Error-Tolerant Quantum State Discrimination}
\label{chap:error_uqsd}
Quantum noise plays a huge role in quantum key distribution. For example, the
B92 protocol under a depolarizing channel has QBER .. (cite, numbers)

When noise is present in the input quantum states, the quantum states can't be
exactly expressed by pure states anymore. While minimum error state discrimination
is still possible, unambiguous state discrimination becomes difficult or even
impossible. (cite)
We will first present the common setting for this scenario, introduce previous
attempts to this scenario, and present our perspectives on the derived problems.

\section{Common setting to discriminate noisy quantum states}
Recall that a quantum state $\Psi$ can be denoted by a state vector $\ket{\psi}$
or a density matrix $\rho$.
The prior probability for each input state $\Psi_{i}$
(expressed by $\ket{\psi_{i}}$ or $\rho_{i}$) is denoted by $p_{i}$, where $i$ ranges from 1 to $k$.
The measurement operators are defined by $\Pi_{i}$, where $i$ ranges from 0 to $k$.
Here $\Pi_{0}$ is designed to match the inconclusive results,
and $\Pi_{1}$ to $\Pi_{k}$ designed to measure the $1$-st to $k$-th state respectively.

Even though a density matrix can express the mixed state, it may provide more
context if it is transformed from a pure state. We assume a depolarizing channel
is applied to the pure state

\section{Fixed rate of inconclusive outcome (FRIO)}
%\cite{origin} \cite{application} \cite{follow-ups}
For discriminating mixed states, a semidefinite programming
problem was proposed \cite{RN37}, and experiments have been performed 
based on this formulation. \cite{RN48, RN221, RN29}

$P_{D}$ denotes the total probability of correct discrimination,
$P_{err}$ denotes the probability of incorrect/erroneous discrimination,
and $P_{inc}$ denotes the probability of inconclusive outcomes.
$\beta$ is a constant real number that is defined to be the target inconclusive probability.
In practice, we use an inequality in place of an equality,
and $\beta$ is the minimum rate of inconclusive outcomes.

We have the following equalities:
\begin{align}
     & P_{D} + P_{err} + P_{inc} = 1 \\
     & \sum_{i=1}^{k} p_{i} = 1
\end{align}

The semidefinite programming problem can be formulated as follows:
\begin{align}
     & \underset{\{\Pi\}}{\text{maximize}} &  & P_{D}=\sum_{i=1}^{k} p_{i} \operatorname{Tr}(\rho_{i} \Pi_{i})              \\
     & \text{subject to}                   &  & \Pi_{i} \succeq 0, \quad i = 0, \ldots, k                                   \\
     &                                     &  & \sum_{i=0}^{k} \Pi_i = I                                                    \\
     &                                     &  & P_{inc}=\sum_{i=1}^{k} p_{i} \operatorname{Tr}(\rho_{i} \Pi_{0}) \geq \beta
\end{align}

One problem regarding to this formulation is that there is a
hyperparameter $\beta$ that is equal to the probability of the
inconclusive result $P_{inc}$. When $\beta$ is set to 0, the problem
degenerates to minimum error discrimination.

However, if we assume that the
mixed states correspond to some pure states, it is likely that we can
accelerate the solving process.

For example...
To view t

\section{Formulation 1}
We are inspired by the type-I error and type-II error in hypothesis testing,
and we include a different set of constraints in the semidefinite programming
problem such that a desired alternative POVM can be obtained. Note that there
is an abundance of discussions about quantum hypothesis testing, which
includes quantum state discrimination, but we will not discuss it for now.
% Assume a set of $k$ input quantum states are expressed in density matrices,
% denoted by $\{ \rho_{1}, \rho_{2}, \dots, \rho_{k}\}$.
The notations are almost identical to FRIO, but our constraints further
consider the error part of the conclusive results.
The formulation also resembles the maximum confidence discrimination scheme
\cite{RN519} and the scheme considering the worst-case a posteriori
probability of detection \cite{kosut2004quantumstatedetectordesign}.

The symbol $O_{i}$ denotes the event that the $i$-th outcome is measured.
The density matrices could be treated without assuming that the states are
related to some pure states. Later, we may assume the mixed states are related
to some pure states or can be expressed by ensembles of pure states.

The semidefinite programming formulation is as follows:
\begin{align}
     & \underset{\{\Pi\}}{\text{maximize}} &  & P_{D}=\sum_{i=1}^{k} p_{i} \operatorname{Tr}(\rho_{i} \Pi_{i})                                                                                                      \\
     & \text{subject to}                   &  & \Pi_{i} \succeq 0, \quad i = 0, \ldots, k                                                                                                                           \\
     &                                     &  & \sum_{i=0}^{k} \Pi_i = I                                                                                                                                            \\
     &                                     &  & P(O_i \mid \rho_i) = \frac{p_i \, \operatorname{Tr}(\rho_i \Pi_i)}{\sum_{j=1}^{L} p_i \, \operatorname{Tr}(\rho_i \Pi_j)} \geq 1 - \alpha_i, \quad i = 1, \ldots, k \\
     &                                     &  & P(\rho_i \mid O_i) = \frac{p_i \, \operatorname{Tr}(\rho_i \Pi_i)}{\sum_{j=1}^{L} p_j \, \operatorname{Tr}(\rho_j \Pi_i)} \geq 1 - \beta_i, \quad i = 1, \ldots, k
\end{align}

\section{Formulation 2}
Since we know the optimal POVM for pure states, we ask if we can find a POVM
such that the resulting probability distribution remains unchanged. First we
need to define the "probability distribution" here and how to evaluate the
distance between two probability distributions.

One of our ideas was to use gradient-based searching methods. We tried to
iterate different hyperparameters and guessed that we would 

(Show gradient-based flow)

\section{Flow 1}

Just copy the SDP form and add the flow chart


- New!! Unbounded/Unconstrained FRIO (E? but no pinc specified) (as a special case)
may benefit in non-orthogonal cases (ex. two states)
每個順便 Discussions
- Other new formulations (many sections)
- Final comparison/summary

Since the pure states are known, we can first 

